\chapter{Введение}
\label{sec:introduction}

Компьютерные сети сегодня — базовая инфраструктура, от надёжности которой зависит работа учреждений и производственных предприятий. С ростом числа подключённых устройств и постепенным внедрением протокола IPv6 наряду с привычным IPv4 растёт и сложность задач, которые администратору сети приходится решать прямо на месте. Вопросы при этом бывают очень конкретными: действительно ли розетка в помещении соединена с распределительным шкафом, какую пропускную способность даёт линия, какие устройства находятся в сегменте и не работают ли на них службы, которых там быть не должно.

Ответ на эти вопросы требует измерения на месте. Именно здесь обычные решения упираются в свои ограничения. Коммерческие ручные тестеры эргономичны, но закрыты: набор их функций задан производителем, и добавить в них собственную методику проверки нельзя. Ноутбук с открытыми инструментами, наоборот, полностью гибок, но в шкафу или под столом с ним работать неудобно, и перед каждым измерением его нужно настраивать вручную. Если же техник попытается подключиться к одному и тому же сегменту двумя интерфейсами, чтобы одновременно слушать трафик и создавать нагрузку, он столкнётся с конфликтом в таблицах маршрутизации операционной системы, который незаметно исказит результаты измерений.

Цель этой бакалаврской работы — спроектировать и реализовать переносное устройство для тестирования и мониторинга компьютерных сетей, которое сочетает открытость программного решения с эргономикой самостоятельного прибора. Прототип построен на одноплатном компьютере Raspberry Pi, оснащён сенсорным дисплеем, собственным управляющим приложением и аккумулятором, поэтому им можно управлять без подключённого компьютера, без физической клавиатуры и без сетевого питания. Он содержит генератор и анализатор трафика, измеряет задержку, пропускную способность, потери и джиттер и позволяет картировать сегмент и проводить разведку безопасности, одинаково для протоколов IPv4 и IPv6.

Ключевой технический элемент конструкции — логическое разделение двух тестовых интерфейсов с помощью технологии Linux Network Namespaces. Это разделение гарантирует, что создаваемая нагрузка не влияет на статистику, захваченную слушающим интерфейсом, и обеспечивается самим ядром операционной системы, а не только аккуратностью оператора при настройке. Второе направление, которому уделено повышенное внимание, — автоматическая настройка адресов в протоколе IPv6, поскольку именно от неё зависит, сможет ли прибор установить связь со станцией, на которой ничего заранее не настроено.

Текст работы разделён на четыре главы. После этого введения следует глава \ref{sec:teorie}, которая излагает теоретические основы: различия протоколов сетевого уровня, адресацию и автоконфигурацию протокола IPv6, генерирование сетевого трафика и измерение его параметров, инструменты мониторинга и аудита трафика и архитектуру пространств имён. Её заключительная часть сравнивает существующие решения и формулирует требования к проектируемому устройству. Глава \ref{sec:realizace} описывает проект прибора с аппаратной и программной стороны, реализацию генератора и анализатора трафика и оценку результатов, проверенных измерениями. Заключение подводит итоги, выделяет собственный вклад автора и намечает направления дальнейшего развития прототипа.
